% New variant section (TRACE emphasis). Reused by main_trace.tex; NOT included in the base main.tex.
\section*{РОЗДІЛ 7. ПОБУДОВА RESEARCHIQ У РАМКАХ TRACE: ІНЖЕНЕРНЕ ОБҐРУНТУВАННЯ ДОВІРЕНОЇ АГЕНТНОЇ СИСТЕМИ}
\label{sec:trace-engineering}

Попередні розділи описали \emph{що} робить ResearchIQ (архітектура, розд.~5) та \emph{як} формується Golden Record (формалізм ER). У цьому розділі ми переходимо до інженерного питання: \emph{на яких підставах} цій агентній системі дозволено діяти в адміністративному контексті, де помилки в метаданих мають реальні наслідки (розд.~2.2.1). Для цього застосовано рамку \textbf{TRACE} --- метрологічно заземлену методологію проєктування довірених агентних AI-систем для операційно-критичних доменів~\cite{Zabolotnii2026TRACE}.

Ключова теза TRACE: \emph{довіра до AI --- не властивість моделі, а властивість системи}. Питання зміщується з «наскільки розумна модель» на: що системі дозволено робити, на якій підставі, під яким наглядом, з якою вимірюваною невизначеністю та чи правильний тип моделі обрано для кожної підзадачі. ResearchIQ розглядається тут як інстанціація TRACE для домену \emph{наукового адміністрування} --- з нижчою критичністю, ніж клінічні, промислові чи судові інстанси~\cite{Zabolotnii2026TRACE}, але з тими самими інженерними вимогами до відтворюваності та підзвітності.

\subsection*{7.1. Шість принципів і чотиришарова архітектура}

TRACE постулює п'ять видимих принципів (акронім T-R-A-C-E) та шостий --- внутрішнє проєктне обмеження: \textbf{T} (Trustworthy / Evidence Traceability), \textbf{R} (Reasoned / Bounded Human Supervision), \textbf{A} (Accountable / Staged Autonomy), \textbf{C} (Context-bound), \textbf{E} (Escalated / Metrological Accountability) та \textbf{Model Parsimony}. Архітектурно вони реалізуються через чотири шари: \textbf{L1} --- детермінований не-навчений core (trust anchor); \textbf{L2} --- stateless-інвентар навчених компонентів, поділений на \textbf{L2a} (класичний ML) і \textbf{L2b} (LLM-валідатори); \textbf{L3} --- окрема stateful policy оркестрації та ескалації; \textbf{L4} --- обмежений людський нагляд із вимірюваним навантаженням.

\subsection*{7.2. Пошарова декомпозиція ResearchIQ}

Принципово, що система ResearchIQ вже побудована за логікою поділу детермінованого ядра й мовної моделі (розд.~3.1.3, 5.1.1): збирання, нормалізація, дедуплікація та злиття реалізуються \emph{детерміновано}, тоді як LLM відповідає за інтерфейс і пояснення. У термінах TRACE (архітектурний зріз --- рис.~\ref{fig:trace-arch}) це відображається так (табл.~\ref{tab:trace-layers}).

\begin{longtable}{@{}p{0.10\textwidth}p{0.40\textwidth}p{0.42\textwidth}@{}}
\caption{Відображення компонентів ResearchIQ на шари TRACE}\label{tab:trace-layers}\\
\toprule
\textbf{Шар} & \textbf{Компонент ResearchIQ} & \textbf{Властивість / роль} \\
\midrule
\endfirsthead
\toprule
\textbf{Шар} & \textbf{Компонент ResearchIQ} & \textbf{Властивість / роль} \\
\midrule
\endhead
\textbf{L1} & \texttt{normalize\_doi}; первинний матчинг за DOI; жорсткі обмеження (різні валідні DOI $\Rightarrow$ заборона auto-merge; $|y_i-y_j|>\Delta_y \Rightarrow$ ручна перевірка); схема провенансу $\mathrm{prov}_a(c)$; версіонування конфігурації $\mathrm{ver}(\Theta)$; audit trail & Trust anchor: не-генеративний, формально перевірюваний; забезпечує traceability (принцип T) \\
\textbf{L2a} & Парний скоринг ER: TF--IDF cosine назв, нормалізований Levenshtein, Jaccard за авторами/venue; логіт $m(r_i,r_j)$ з фіксованими або навченими вагами & Структуроване передбачення з калібровними ймовірностями, низькою затримкою, відтворюваністю \\
\textbf{L2b} & Інтерпретація NL-запитів; генерація пояснень і чернеток звітів; семантична звірка вільнотекстових нотаток; перевірка галюцинацій у згенерованих резюме & Застосовується \emph{лише} там, де потрібне лінгвістичне розуміння free-form тексту \\
\textbf{L3} & Двопорогова політика $\theta_{\mathrm{match}}/\theta_{\mathrm{review}}$ (auto-merge / flag); fail-fast на конфлікті високонадійних джерел; тригери HITL; re-invocation з розширеним контекстом & Stateful policy оркестрації та ескалації над інвентарем L2 \\
\textbf{L4} & HITL-контур: попередній перегляд перед фінальним звітом; маркування записів низької впевненості; override будь-якого авто-рішення; аудит-лог правок & Фінальна підзвітність; reviewer --- адміністратор/відповідальний кафедри \\
\bottomrule
\end{longtable}

\subsection*{7.3. Явне рішення L2a проти L2b (Model Parsimony)}

Центральна інженерна вимога TRACE --- свідомий вибір типу навченого компонента за \emph{fitness for purpose}, а не за LLM-презумпцією. ResearchIQ задовольняє її конструктивно: ядро метрик і трансформацій (дедуплікація, скоринг, агрегація показників) \emph{не покладається на LLM}, оскільки для structured prediction детерміновані методи та класичний ML дають кращу калібровку, нижчу затримку й перевірювану відтворюваність. LLM (L2b) застосовується точково --- для NL-запитів, пояснень і верифікації вільного тексту, тобто там, де лінгвістичне розуміння справді необхідне. Це пряме уникнення антипатерну «LLM-by-default»: спроба виконувати дедуплікацію чи розрахунок H-індексу через LLM дала б гірший результат і непрозору похибку.

\subsection*{7.4. L3: політика ескалації окремо від інвентаря моделей}

ResearchIQ містить діючу ескалаційну логіку: два пороги $\theta_{\mathrm{match}}=0.90$, $\theta_{\mathrm{review}}=0.70$ розмежовують автоматичне злиття та передачу на ручну перевірку, а fail-fast при двох різних валідних DOI запобігає катастрофічним злиттям. Важливо, що ескалація спирається не лише на низьку впевненість: конфлікт між високонадійними джерелами ескалюється навіть за високої впевненості (відповідає принципу TRACE «escalate on joint risk $\wedge$ confidence»).

\textbf{Обмеження поточної реалізації (evidence gap).} У версії розд.~5.1.3 послідовність викликів інструментів визначає сам LLM-агент. За TRACE це змішує concern рівня L2 (вибір моделі) з concern рівня L3 (invocation policy). Рекомендація: винести L3 у детермінований/правиловий контролер над інвентарем L2, лишивши LLM роль валідатора в його межах; зміна L2a$\leftrightarrow$L2b тоді не вимагатиме переписування L3.

\subsection*{7.5. L4: людський нагляд із вимірюваним навантаженням}

Принцип TRACE вимагає, щоб людський нагляд був архітектурним шаром із чітким reviewer-ом, override-правами та \emph{вимірюваним} навантаженням, а не косметичним «human-in-the-loop». ResearchIQ реалізує контур попереднього перегляду, маркування низької впевненості та повного override (розд.~5.3.4), а також аудит-лог. Для відповідності TRACE цей шар має супроводжуватися метриками навантаження (RBI) та частки перевизначень (OvR): значення $\mathrm{OvR}\approx 0$ свідчило б про «rubber-stamp», а не про якісний нагляд.

\subsection*{7.6. Підмножина trust-метрик для ResearchIQ}

TRACE специфікує 17 метрик; для системи слід обирати релевантну підмножину, а не вимірювати все. Табл.~\ref{tab:trace-metrics} фіксує мінімальний протокол для ResearchIQ та узгоджує його з протоколом оцінювання з розд.~5.2.2. Більшість значень --- це \emph{план вимірювання} (узгоджено з \textit{Evidence gap} розд.~5.2), окрім архітектурно обґрунтованого CPR (розд.~7.7).

\begin{longtable}{@{}p{0.10\textwidth}p{0.20\textwidth}p{0.62\textwidth}@{}}
\caption{Підмножина trust-метрик TRACE для ResearchIQ}\label{tab:trace-metrics}\\
\toprule
\textbf{Шар} & \textbf{Метрика} & \textbf{Операціоналізація в ResearchIQ} \\
\midrule
\endfirsthead
\toprule
\textbf{Шар} & \textbf{Метрика} & \textbf{Операціоналізація в ResearchIQ} \\
\midrule
\endhead
L1 & RCR, UTC & Частка записів, розв'язаних детермінованими правилами (DOI / Title+Year) без LLM; частка merge-рішень із задокументованим rationale та провенансом (схема $\mathrm{prov}_a(c)$ це вже підтримує) \\
L2 & CRP, CE & Релевантність контексту, переданого LLM; \textbf{калібрування merge-впевненості} $\gamma_a$ через ECE/Brier --- наразі \textit{Evidence gap} (ECE не виміряно) \\
L3 & EP, TCC & Частка flag-for-review, що справді потребували людини; обчислювальна вартість шляху на запит \\
L4 & RBI, OvR & Час reviewer-а на звіт/запис; частка авто-злиттів, перевизначених людиною \\
Cross & ETC, OSI & Повнота evidence-trail до першоджерел (провенанс уже фіксується); дрейф ключових метрик у часі \\
Economy & \textbf{CPR} & Computational Parsimony Ratio (обов'язково, розд.~7.7) \\
\bottomrule
\end{longtable}

Бенчмарк ResearchIQ вже звітує довірчі інтервали біля кожного середнього значення; у термінах TRACE/GUM їх доцільно переоформити як розширену невизначеність із коефіцієнтом покриття~$k$, замкнувши метрологічний контур.

\subsection*{7.7. CPR: обчислення з наявного бенчмарку}

CPR визначається як відношення вартості найекономічнішої моделі, що адекватно вирішує задачу, до вартості розгорнутої моделі; $\mathrm{CPR}\approx 1$ означає оптимальний вибір, $\mathrm{CPR}\ll 1$ --- архітектурний overhead і ймовірний LLM-by-default bias~\cite{Zabolotnii2026TRACE}. Для \emph{ядра консолідації даних} ResearchIQ ця діагностика є майже безкоштовною: бенчмарк уже розділяє compute, фіксуючи час двигуна дедуплікації та злиття ($\approx 1$~с на факультет 50 НПП, Додаток~В, табл.~\ref{tab:efficiency-50}) \emph{окремо від} мережевих і LLM-витрат. Саме це детерміноване ядро автоматизує $\approx 96\%$ модельованих витрат ручної праці без жодного звернення до LLM. Оскільки це ядро реалізоване детерміновано (L1+L2a), для нього $\mathrm{cost}_{\mathrm{minimal}}\approx\mathrm{cost}_{\mathrm{deployed}}$, тобто $\mathrm{CPR}\approx 1$: система \emph{за побудовою} не витрачає LLM там, де достатньо детермінованих методів.

\textbf{Уточнення (evidence gap).} Повний CPR має враховувати й вартість інтерфейсного (L2b) шляху; поточний бенчмарк вимірює лише детерміноване ядро. Тому ми стверджуємо $\mathrm{CPR}\approx 1$ \emph{для ядра консолідації} та фіксуємо вартість інтерфейсного шляху як наступний крок вимірювання (узгоджено з \textit{Evidence gap} щодо cost-per-query, розд.~6.1.4). У регульованому розгортанні CPR є auditable-артефактом для обґрунтування вибору моделі (EU AI Act, ст.~15).

\subsection*{7.8. Відображення на governance-рамки}

TRACE операціоналізує регуляторні обов'язки як інженерні артефакти. Для ResearchIQ найрелевантніші відповідності: EU AI Act ст.~12--13 (logging/transparency) $\rightarrow$ ETC + UTC + evidence chain; ст.~14 (human oversight) $\rightarrow$ L4 + RBI + OvR + ABC; ст.~15 (accuracy/robustness) $\rightarrow$ метрики L2 + CE + \emph{CPR як efficiency-артефакт}; ст.~9 (risk management) $\rightarrow$ L3-policy + L4 + OSI. Метрологічну основу всіх метрик задають GUM/VIM/ISO~17025, а систему загалом можна позиціонувати як реалізацію AIMS за ISO/IEC~42001~\cite{Zabolotnii2026TRACE}. Вимоги GDPR та EDPB Opinion 28/2024 розглянуто в розд.~6 (ризики), тож тут вони не дублюються. Зауважимо, що наукове адміністрування ймовірно належить до \emph{limited-risk}, а не high-risk за AI Act, проте інженерна дисципліна TRACE застосовна незалежно від класу ризику.

\subsection*{7.9. Антипатерни: чого ResearchIQ уникає і на що зважати}

За рамкою TRACE ResearchIQ \emph{уникає} двох поширених антипатернів: «LLM-by-default» (ядро ER детерміноване, а не мовно-модельне) та «learned L1» (правила нормалізації й жорсткі обмеження задані явно, а не «вивчені» моделлю). Водночас слід зважати на три ризики: (i) \emph{L3/L2b-конфляція} --- наразі оркестрацію виконує LLM, що TRACE рекомендує відокремити в детермінований L3; (ii) \emph{косметичний L4} --- нагляд має підтверджуватися метрикою $\mathrm{OvR}>0$, а не лише наявністю кнопки override; (iii) \emph{невиміряна калібровка} --- бал впевненості $\gamma_a$ слід калібрувати (ECE), інакше пороги ескалації L3 спираються на неперевірену впевненість. Усунення цих трьох пунктів переводить ResearchIQ з «системи, що відповідає духу TRACE» у «систему з вимірюваною відповідністю TRACE».
